\section{Un mismo modelo, dos canales de aprendizaje}

Esta clase parte de una pregunta que parece sencilla pero es fundamental: cuando un modelo de lenguaje recibe nueva información, ¿produce el mismo tipo de conocimiento accesible \textbf{escribiéndola en los parámetros} o \textbf{dejándola en el contexto}? Ambas rutas permiten que el modelo se comporte como si hubiera «aprendido» en la tarea actual, pero los experimentos controlados de Lampinen y colaboradores muestran que, en pruebas que requieren reorganizar relaciones—como reversal, syllogism, codebook o alternative goal—ambas exhiben patrones de generalización significativamente distintos.

\subsection{El objeto de estudio no es solo el Transformer, también incluye al propio sistema de aprendizaje}

La clase sitúa primero la pregunta en un marco común entre la inteligencia natural y la artificial. No se trata de inferir mecanismos cerebrales a partir de un experimento con Transformers, sino de aprovechar la controlabilidad de los sistemas artificiales: los investigadores pueden introducir \textbf{la misma nueva información} tanto en el contexto como en los parámetros y observar cómo cambian las predicciones posteriores. En humanos, es difícil controlar con precisión hacia qué tipo de sistema de memoria ha entrado una experiencia; en los modelos, esta intervención puede repetirse, eliminarse (ablation) y cuantificarse.

\lecturefigure{slide-002.jpg}{Inferir principios computacionales del sistema de aprendizaje a partir de la generalización del modelo de lenguaje}{Diapositiva oficial de Stanford CS25 V6, clase 06, pág. 2}

\teachervoice{00:01:06--00:02:18，el docente enfatiza que se trata de un método de ciencia cognitiva: el valor de los modelos no reside únicamente en su desempeño como objetos de prueba, sino también en su capacidad de ser sistemas de aprendizaje sujetos a intervenciones controladas. Las analogías con la inteligencia natural deben establecerse después de estos experimentos con modelos y no consistir en asignar nombres de regiones cerebrales directamente a componentes del modelo.}

\begin{importantbox}{Pregunta central de esta clase}
Dada la misma experiencia $D$, si un modelo aprende $D$ mediante actualización de parámetros, ¿podría realizar inferencias igualmente flexibles sobre las \textbf{relaciones no explícitamente escritas} que si se introduce $D$ directamente en el contexto? Si no es así, ¿provienen las diferencias de la capacidad, la optimización, la forma de acceso o la disponibilidad de la representación del conocimiento?
\end{importantbox}

\subsection{Aprendizaje paramétrico y aprendizaje in-contexto}

El aprendizaje paramétrico comprime grandes cantidades de muestras en los pesos mediante la minimización de una función de pérdida durante el entrenamiento; el aprendizaje in-contexto (in-context learning, ICL) no actualiza los pesos, sino que utiliza hechos, ejemplos y documentos dentro del prompt durante una sola pasada de inferencia hacia adelante. Ambos pueden modificar la salida, pero corresponden a diferentes ciclos de vida de la información: el primero es persistente y amortizable, aunque la experiencia quede altamente comprimida; el segundo es efímero, ocupa la ventana de contexto, pero conserva más detalles locales.

\lecturefigure{slide-003.jpg}{Las dos rutas de aprendizaje en modelos de lenguaje: optimización paramétrica y uso del contexto}{Diapositiva oficial de Stanford CS25 V6, clase 06, pág. 3}

El aprendizaje paramétrico se puede expresar como
$$
\theta' = \arg\min_{\theta}\sum_{(x,y)\in D}\mathcal{L}\!\left(f_{\theta}(x),y\right),
$$
donde:
\begin{itemize}
  \item $D$ es el nuevo conjunto de datos de entrenamiento;
  \item $f_{\theta}$ es un modelo de lenguaje con parámetros $\theta$;
  \item $\mathcal{L}$ es la pérdida siguiente al token o la tarea;
  \item $\theta'$ son los nuevos parámetros que retienen la información por más tiempo tras absorber los datos.
\end{itemize}

ICL mantiene los parámetros invariables:
$$
p(y\mid x,D)=f_{\theta}\!\left(y\mid \operatorname{concat}(D,x)\right).
$$
En este caso, $D$ no se comprime dentro de $\theta$, sino que participa directamente en la predicción como parte del grafo de cálculo actual. Por ello, «qué sabe el modelo» debe descomponerse en dos preguntas: si la información existe y si puede reorganizarse y activarse bajo el objetivo actual.

\begin{knowledgebox}{Glosario: ¿en qué difieren realmente las dos formas de «aprender»?}
\begin{tabularx}{\textwidth}{L{0.22\textwidth}X X}
\toprule
Dimensión & Aprendizaje paramétrico & Aprendizaje in-contexto \\
\midrule
Ubicación de la información & Pesos y representaciones internas & Secuencia actual de tokens y sus activaciones \\
Duración & Se conserva a través de múltiples solicitudes & Generalmente válido solo en el contexto actual \\
Alcance de integración & Integración estadística a través de millones de muestras de entrenamiento & Limitado por la ventana de contexto y la calidad de la recuperación \\
Conservación de detalles & Comprimidos mediante optimización, pueden quedar ligados al formato original de la tarea & La experiencia original se conserva más íntegra y puede releerse bajo nuevos objetivos \\
Costo principal & Entrenamiento, actualización y gestión de versiones & Tokens de inferencia, caché KV, recuperación y latencia \\
\bottomrule
\end{tabularx}
\end{knowledgebox}

\subsection{Experimentos controlados: cambiar solo la fuente de la información}

El método clave de esta clase no es comparar dos modelos distintos, sino fijar el modelo, los datos y las pruebas, variando únicamente el canal de entrada de la información. Los investigadores pueden construir entidades y relaciones que el modelo nunca vio durante el preentrenamiento, introducir todo el conjunto de nuevos datos en el contexto, o realizar un ajuste fino con los mismos datos, para luego probar por separado hechos directos, relaciones inversas, razonamiento compuesto y estructuras latentes. Este diseño separa la «fuente de la información» del tamaño del modelo, el conocimiento del mundo y los corpus de entrenamiento.

\lecturefigure{slide-004.jpg}{Experimento controlado que inyecta la misma información novel en el contexto o en los parámetros}{Diapositiva oficial de Stanford CS25 V6, clase 06, pág. 4}

\begin{lstlisting}[language=Python,caption={Pseudocódigo para comparación controlada con los mismos datos pero distintos canales de información}]
dataset = build_novel_relations()
tests = build_forward_and_latent_tests(dataset)

icl_model = frozen_model
icl_scores = evaluate(icl_model, context=dataset, tests=tests)

ft_model = finetune(copy(frozen_model), dataset)
ft_scores = evaluate(ft_model, context=None, tests=tests)

compare(icl_scores, ft_scores, by=["forward", "reversal", "composition"])
\end{lstlisting}

\lecturefigure{slide-005.jpg}{Ruta de la clase: medir primero las diferencias, luego comparar tres esquemas de puente y finalmente discutir la inteligencia natural}{Diapositiva oficial de Stanford CS25 V6, clase 06, pág. 5}

\begin{warningbox}{No equiparar directamente «saber responder» con «haber formado el mismo tipo de conocimiento»}
El hecho de que dos sistemas respondan correctamente en la distribución de entrenamiento o en preguntas explícitas no implica que su conocimiento interno tenga la misma capacidad de reorganización. Lo verdaderamente distintivo es si el rendimiento se mantiene cuando se invierte la dirección de la pregunta, se cambia el objetivo, se altera el idioma, se añade una relación adicional o cuando se exige llamar a conclusiones que solo estaban implícitas durante el entrenamiento.
\end{warningbox}

\subsection{Resumen de la sección}

En este capítulo se establecen las variables causales para toda la clase: la misma información entra al modelo ya sea por los parámetros o por el contexto. El aprendizaje paramétrico es bueno para integrar estadísticamente grandes cantidades de experiencias a largo plazo; ICL permite que pocas experiencias participen directamente en el cálculo actual con una riqueza de detalles mayor. A continuación no se trata de averiguar qué camino es «más inteligente», sino de identificar qué tipo de pruebas expone las diferencias entre ellas en cuanto a la disponibilidad del conocimiento.


